% ============================================================================
% Fichier  : partie1/P01_C02_grandes_defaillances.tex
% Titre    : Grandes Défaillances et Naissance de l'Audit des SI
% Auteur   : MINKA MI NGUIDJOÏ Thierry Emmanuel
% Version  : 1.0 -- novembre 2026
% ============================================================================
% Chapitre-cas. Sources vérifiées le 01/11/2026 :
% - Enron : 8-K du 8/11/2001 (retraitement, -586 M$ de résultats 1997-2001
%   S1, +2,59 Md$ de dettes) ; faillite le 2/12/2001 ; rapport Powers du
%   1/2/2002 ; condamnation d'Andersen en juin 2002 pour entrave (destruction
%   de documents), annulée par la Cour suprême le 31/05/2005 ; SOX du
%   30/07/2002, section 404 (15 U.S.C. 7262, texte lu sur law.cornell.edu).
% - Société Générale : document de la banque « Les 10 points clés » (juin
%   2016 : 18-20/01/2008, positions de près de 50 Md EUR, perte 4,9 Md EUR =
%   gain dissimulé 1,4 Md - pertes de débouclage 6,3 Md ; limite de l'équipe
%   125 M EUR ; « 75 alertes » ; contrôles dispersés ; plan Fighting Back ;
%   décisions du 5/10/2010, 24/10/2012, 19/03/2014) ; Commission bancaire,
%   blâme et 4 M EUR, 4/07/2008 (presse concordante).
% - Equifax : GAO-18-559 (texte lu) : vulnérabilité Apache Struts publiée
%   le 8/03/2017, scan des attaquants le 10/03, accès à partir du 13/05,
%   découverte le 29/07, environ 76 jours, 3 + 48 bases, environ 9 000
%   requêtes, certificat expiré environ 10 mois avant, liste de diffusion
%   périmée, scan inefficace ; au moins 145,5 M de personnes ; règlement
%   jusqu'à 700 M$ (CFPB, 22/07/2019). Le rapport de la Chambre (déc. 2018)
%   donne 19 mois pour le certificat : divergence signalée dans le texte.
% - Wirecard : rapport KPMG remis le 28/04/2020 (communiqué de la société) ;
%   refus d'EY le 18/06/2020 ; 1,9 Md EUR ; déclaration du 22/06 ;
%   insolvabilité le 25/06/2020 ; interdiction des ventes à découvert par la
%   BaFin en février 2019 ; commission d'enquête du Bundestag (2020).
% Knight Capital (2012) et TSB (2018) sont traités dans le cours GCI.
% Corrigés : corriges/Corr_C02.tex.
% ============================================================================

\chapter[Grandes Défaillances et Naissance de l'Audit des SI]{Grandes Défaillances et Naissance de l'Audit des SI}
\label{chap:defaillances}
\chaptermark{Grandes défaillances}

\epigraph{%
    « Chaque catastrophe a eu ses contrôles.\\
    Ils existaient sur le papier~; personne n'a vérifié qu'ils fonctionnaient. »%
}{Leçon commune des quatre affaires de ce chapitre}

\niveauChapitre{Fondamental}{%
    Maîtrise du Chapitre~\ref{chap:audit-auditologie} (critère, preuve,
    constatation~; assurance~; trois lignes)%
}

\begin{boxobjectifs}
\begin{itemize}
    \item Lire une affaire réelle avec méthode~: séparer les faits établis,
          leurs sources et leur interprétation.
    \item Distinguer une source primaire d'une source secondaire, et une
          source neutre d'une source partie au litige.
    \item Analyser quatre affaires documentées (Enron, Société Générale,
          Equifax, Wirecard)~: le mécanisme, le contrôle qui a manqué,
          l'audit qui ne l'a pas vu, la réforme qui a suivi.
    \item Expliquer pourquoi la loi Sarbanes-Oxley a fait des contrôles
          informatiques un objet d'audit, et ce que dit sa section~404.
    \item Reconnaître cinq schémas de défaillance qui reviennent d'une
          affaire à l'autre, et les retrouver dans le cas LiZbiZ.
    \item Relier chaque schéma à la réponse \NEK{}~: une alerte fermée par
          une déclaration est un \NEK{} déguisé en \OUI{}.
\end{itemize}
\end{boxobjectifs}

\bigskip

% ============================================================================
\section*{Question directrice : pourquoi étudier des catastrophes~?}
% ============================================================================

Un ingénieur en génie civil étudie les ponts qui se sont effondrés~; un
ingénieur aéronautique, les rapports d'accident. Ils n'y cherchent pas
des coupables, mais des mécanismes~: ce qui a cédé, pourquoi les
protections n'ont pas joué, ce que l'on a changé ensuite. L'auditeur fait
de même. Les grandes affaires de l'histoire récente ne sont pas des
anecdotes~: chacune a laissé une règle, une norme ou une loi que l'auditeur
applique aujourd'hui.

Ce chapitre en étudie quatre, choisies pour couvrir des terrains
différents~: un montage comptable (Enron, 2001), une fraude interne dans une
salle de marchés (Société Générale, 2008), une cyberattaque (Equifax,
2017), une fraude comptable dans une société de paiement (Wirecard, 2020).
Deux autres affaires, Knight Capital (2012) et la migration informatique
de la banque TSB (2018), sont étudiées dans le cours MSN-2S3-MSI03 sous
l'angle de la gouvernance.

Le chapitre commence par une méthode de lecture
(section~\ref{sec:c02-methode}), raconte les quatre affaires
(sections~\ref{sec:c02-enron} à~\ref{sec:c02-wirecard}), puis en tire les
schémas communs (section~\ref{sec:c02-schemas}).

\separateur

% ============================================================================
\section{Lire une Affaire avec Méthode}
\label{sec:c02-methode}
% ============================================================================

\subsection{Six questions}

Une affaire se lit comme une mission d'audit~: avec une grille, et en
séparant ce qui est établi de ce que l'on en pense. Le cours utilise six
questions.

\begin{tableaucours}{cL{4.2cm}L{7.4cm}}{%
    \tableauHeader{N°} &
    \tableauHeader{Question} &
    \tableauHeader{Ce qu'on cherche}
}
1 & Que s'est-il passé~? & Les faits datés, chiffrés, sourcés \\
2 & Quel était le mécanisme~? & Comment la perte, la fraude ou la fuite
s'est produite techniquement \\
3 & Quel contrôle a manqué~? & Le contrôle absent, mal conçu, ou conçu mais
inopérant \\
4 & Quel signal a été ignoré~? & L'alerte, le doute, la question qui aurait
permis de voir \\
5 & Pourquoi l'audit n'a-t-il pas vu~? & Indépendance, preuve, périmètre,
méthode \\
6 & Qu'est-ce qui a changé ensuite~? & Loi, norme, pratique de la
profession \\
\end{tableaucours}
\vspace{-0.8em}
\begin{center}{\small\itshape Tableau~2.1. La grille de lecture d'une
affaire}\end{center}

\subsection{La critique des sources}

Une affaire célèbre est racontée par des centaines de sources, qui ne
disent pas toutes la même chose. L'auditeur applique aux sources la même
discipline qu'aux preuves (Chapitre~\ref{chap:preuve}).

\begin{boxdef}
\termecle{Source primaire}\index{source primaire}~: document produit par
l'acteur ou l'autorité qui a établi le fait~: décision de justice, rapport
d'une autorité de contrôle, déclaration réglementaire de l'entreprise,
texte de loi.

\medskip
\termecle{Source secondaire}\index{source secondaire}~: document qui
rapporte ou commente une source primaire~: article de presse, ouvrage,
encyclopédie en ligne.

\medskip
\termecle{Source partie}\index{source partie}~: source primaire produite par
un acteur qui a un intérêt dans l'affaire (l'entreprise, l'accusé). Elle
établit ce que cet acteur affirme~; elle ne suffit pas à établir que c'est
vrai.
\end{boxdef}

Dans ce chapitre, chaque fait est rattaché à une source primaire, citée en
bibliographie. Quand une source est partie au litige, le texte le dit
(«~selon la banque~»). Quand deux sources divergent, le texte le signale au
lieu de choisir en silence.

\begin{boxalert}
\textbf{Violation typique.} Reprendre dans un rapport ou un mémoire un
chiffre lu dans un article, sans remonter à la source, puis le voir repris
par d'autres.

\smallskip
\textbf{Règle~:} un chiffre n'est cité que si l'on a lu sa source
primaire. À défaut, on écrit «~selon la presse~» ou «~à vérifier~», comme
on cote \NEK{} une réponse que l'on ne peut pas prouver.
\end{boxalert}

% ============================================================================
\section{Enron (2001) : le Montage que Personne ne Lisait}
\label{sec:c02-enron}
% ============================================================================

\subsection{Les faits}

Enron était, à la fin des années~1990, l'une des plus grandes entreprises
américaines, active dans le négoce de l'énergie. Le 8~novembre 2001, elle
déclare à l'autorité boursière américaine (\sigle{SEC}) qu'elle retraite
ses comptes de 1997 au premier semestre 2001~: les résultats publiés sont
réduits d'environ 586~millions de dollars, et ses dettes augmentées
d'environ 2,59~milliards, principalement parce que des sociétés partenaires
auraient dû être consolidées~\cite{enron2001restatement}. Le 2~décembre
2001, Enron se place sous la protection de la loi sur les faillites.

Le conseil d'administration avait créé, fin octobre~2001, un comité
spécial d'enquête, présidé par William Powers. Son rapport, rendu public le
1\textsuperscript{er}~février 2002, décrit les opérations conclues avec ces
sociétés partenaires et les défaillances qui les ont
permises~\cite{powers2002enron}.

\subsection{Le mécanisme}

Les sociétés partenaires étaient des \termecle{entités ad
hoc}\index{entité ad hoc}~: des structures créées pour une opération
précise, et que les règles comptables de l'époque permettaient, sous
conditions, de ne pas consolider. Enron leur cédait des actifs ou concluait
avec elles des couvertures, enregistrait des résultats, et laissait les
dettes hors de son bilan. Plusieurs de ces entités étaient gérées par des
dirigeants d'Enron, dont le directeur financier~: l'entreprise traitait,
de fait, avec elle-même.

\subsection{Les contrôles qui ont manqué}

Le rapport Powers met en cause plusieurs niveaux~\cite{powers2002enron}.
\begin{itemize}
    \item \textbf{La gouvernance}~: le conseil a autorisé le directeur
          financier à diriger des entités qui traitaient avec Enron, ce qui
          supposait de déroger au code de conduite, sans mettre en place une
          surveillance à la hauteur du conflit d'intérêts.
    \item \textbf{Le contrôle des opérations}~: les contrôles prévus sur ces
          opérations (revues, approbations) n'ont pas fonctionné comme
          prévu, et n'ont pas été vérifiés.
    \item \textbf{L'information du conseil}~: le comité d'audit recevait des
          présentations trop générales pour comprendre des montages d'une
          grande complexité.
\end{itemize}

\subsection{L'audit qui n'a pas vu}

Le commissaire aux comptes d'Enron était Arthur Andersen, alors l'un des
cinq grands cabinets mondiaux. Le cabinet a approuvé la comptabilisation de
plusieurs de ces montages. En juin~2002, il a été condamné pour entrave à
la justice, pour avoir détruit des documents liés à Enron~; condamnation
annulée en mai~2005 par la Cour suprême des États-Unis pour un défaut des
instructions données au jury. Entre-temps, le cabinet avait cessé ses
activités d'audit et disparu. La faute d'audit la plus souvent analysée est
celle de l'indépendance~: un cabinet qui tirait aussi d'Enron d'importants
honoraires de conseil était mal placé pour contredire son client.

\subsection{Ce qui a changé~: la loi Sarbanes-Oxley}

Le 30~juillet 2002, les États-Unis adoptent la loi
Sarbanes-Oxley~\cite{sox2002}. Trois de ses dispositions concernent
directement l'audit.
\begin{itemize}
    \item Elle crée un organe de supervision des auditeurs des sociétés
          cotées, le \sigle{PCAOB}, qui fixe leurs normes et les inspecte.
    \item Elle interdit à l'auditeur de fournir à son client certains
          services autres que l'audit.
    \item Sa section~404 impose un rapport sur le contrôle interne.
\end{itemize}

\begin{boxjuris}
\textbf{Sarbanes-Oxley Act, section~404 (a) et (b)} (15 U.S.C.~7262,
traduction résumée)~: le rapport annuel de chaque société cotée contient un
rapport sur le contrôle interne, qui (1) affirme la responsabilité de la
direction dans la mise en place et le maintien d'un dispositif de contrôle
interne et de procédures adaptés pour l'information financière, et (2)
évalue, à la clôture de l'exercice, l'efficacité de ce dispositif. Le
cabinet qui audite les comptes atteste l'évaluation faite par la direction
et en rend compte.

\medskip
\textbf{Portée pour l'audit des SI.} L'information financière est produite
par des applications. Évaluer le contrôle interne sur l'information
financière, c'est donc évaluer les contrôles informatiques~: qui peut
modifier un programme, qui peut accéder aux données, comment
l'exploitation est surveillée. Les \emph{contrôles généraux
informatiques} sont devenus, depuis, un passage obligé de l'audit
(Chapitre~\ref{chap:controle-interne}).
\end{boxjuris}

\begin{boiteRemarque}{Et hors des États-Unis~?}
La loi Sarbanes-Oxley ne s'applique qu'aux sociétés cotées aux États-Unis.
Mais elle a inspiré, dans de nombreux pays, des règles sur le comité
d'audit, l'indépendance des commissaires aux comptes et le rapport sur le
contrôle interne. Le droit OHADA et la réglementation bancaire de la CEMAC
en portent la marque, à leur manière (Chapitre~\ref{chap:droit}).
\end{boiteRemarque}

% ============================================================================
\section{Société Générale (2008) : les Alertes Éteintes}
\label{sec:c02-socgen}
% ============================================================================

\subsection{Les faits}

Les faits ci-dessous sont ceux qu'établit la banque, partie au litige,
dans un document publié en 2016~\cite{sg2016kerviel}~; les décisions de
justice qu'il cite ont confirmé la culpabilité de l'auteur de la fraude.

\begin{tableaucours}{L{2.8cm}L{9.4cm}}{%
    \tableauHeader{Date} &
    \tableauHeader{Fait (selon la banque)}
}
2007 & Un trader de l'équipe dite \emph{Delta One} prend des positions non
autorisées qui atteignent environ 30~milliards d'euros~; à la fin de
l'année, elles dégagent un gain latent de 1,4~milliard, dissimulé \\
Janvier 2008 & Nouvelles positions non autorisées, de près de 50~milliards
d'euros, plus que les fonds propres de la banque \\
18 janvier 2008 & Une alerte conduit à une enquête interne~; les 19 et
20~janvier, les positions dissimulées sont découvertes \\
21-23 janvier & Débouclage des positions, en trois jours \\
24 janvier 2008 & Annonce publique~: perte de 4,9~milliards d'euros (pertes
de débouclage de 6,3~milliards, moins le gain dissimulé de 1,4) \\
4 juillet 2008 & La Commission bancaire inflige à la banque un blâme et une
amende de 4~millions d'euros pour des carences de son contrôle interne \\
2010-2014 & Le trader est condamné, notamment pour faux et usage de faux et
pour introduction frauduleuse de données dans un système de traitement
automatisé (jugement du 5~octobre 2010, arrêts du 24~octobre 2012 et du
19~mars 2014) \\
\end{tableaucours}
\vspace{-0.8em}
\begin{center}{\small\itshape Tableau~2.2. L'affaire de la Société Générale
(d'après~\cite{sg2016kerviel} et la décision de la Commission bancaire)}\end{center}

La limite de risque de l'équipe, huit traders, était de 125~millions
d'euros en fin de journée~\cite{sg2016kerviel}. Les positions découvertes
la dépassaient d'environ quatre cents fois.

\subsection{Le mécanisme}

Pour masquer ses positions, le trader saisissait des \emph{opérations
fictives}~: des opérations de sens inverse, qui faisaient apparaître dans
les systèmes une position nette faible. Selon la banque, il en a enregistré
plusieurs centaines, qu'il annulait et remplaçait avant qu'elles ne
déclenchent les contrôles de confirmation~\cite{sg2016kerviel}. Il avait
été assistant pendant plus de deux ans dans les fonctions de suivi des
opérations, et connaissait de l'intérieur les contrôles qu'il déjouait.

\subsection{Le contrôle qui a manqué~: les alertes}

C'est ici que l'affaire intéresse le plus l'auditeur. Les contrôles
\emph{ont} réagi. La banque évoque elle-même \textbf{75~alertes} sur la
période. Chacune a été fermée~: le trader fournissait une explication, et
parfois un faux, par exemple de faux courriels (la banque en mentionne
sept), et les services de contrôle acceptaient l'explication. La banque
reconnaît que ses contrôles étaient «~dispersés au sein de plusieurs
directions différentes sans vision transversale~»~\cite{sg2016kerviel}.

\begin{boxscenario}{Une alerte, trois réponses}
\textbf{Situation.} Un contrôle de second niveau signale une opération sans
contrepartie confirmée. Le trader répond par courriel qu'il s'agit d'une
erreur de saisie, corrigée le jour même.

\medskip
\textbf{Réponse 1 (acceptée en 2007).} L'explication est plausible~:
l'alerte est close, cotée de fait \OUI{}.

\textbf{Réponse 2.} L'explication n'est pas prouvée~: l'alerte reste
ouverte, cotée \NEK{}, jusqu'à confirmation par la contrepartie.

\textbf{Réponse 3.} La confirmation est demandée directement à la
contrepartie, sans passer par le trader~: la réponse établit l'opération,
ou révèle qu'elle n'existe pas.

\medskip
\textbf{Analyse} (confiance élevée). La réponse~1 transforme une
déclaration en preuve~: c'est exactement l'erreur de la première mission
de LiZbiZ. Une alerte fermée par une explication non vérifiée est un
\NEK{} déguisé en \OUI{}. La réponse~3 est celle de l'auditeur~: elle
obtient la preuve auprès d'une source indépendante de la personne
contrôlée.
\end{boxscenario}

\subsection{Ce qui a changé}

Selon la banque, le plan de remédiation a porté sur le suivi des risques en
montant nominal (et non seulement en position nette), le suivi des
annulations d'opérations, le suivi de la prise des congés, et la création
d'un département transversal de supervision de la sécurité des
opérations~\cite{sg2016kerviel}. Chacune de ces mesures répond à un
mécanisme de la fraude~:
\begin{itemize}
    \item le \emph{nominal} rend visibles des positions que le net
          masquait~;
    \item le suivi des \emph{annulations} détecte les opérations fictives
          remplacées en chaîne~;
    \item les \emph{congés} obligatoires interrompent une fraude qui exige
          une présence continue pour être entretenue~;
    \item la \emph{centralisation} des alertes permet de voir qu'une même
          personne en génère beaucoup.
\end{itemize}

\begin{boxmetier}{Contrôleur interne}
\textbf{Ce que l'affaire lui apprend}
\begin{enumerate}
    \item Une alerte n'est fermée que sur preuve, jamais sur explication.
    \item Les alertes se centralisent et se comptent par personne~: dix
          alertes « expliquées » chez le même opérateur forment un signal.
    \item Un ancien du suivi des opérations qui passe à la prise de risque
          connaît les contrôles~: c'est un facteur de risque à suivre, pas
          une garantie de rigueur.
    \item Les congés obligatoires sont un contrôle, et leur respect se
          vérifie.
\end{enumerate}
\end{boxmetier}

% ============================================================================
\section{Equifax (2017) : le Correctif qui n'est pas Arrivé}
\label{sec:c02-equifax}
% ============================================================================

\subsection{Les faits}

Equifax est l'une des trois grandes agences américaines d'évaluation du
crédit des particuliers. Elle détient des données personnelles sur la
majorité des adultes du pays. Le rapport de l'office de contrôle du
Congrès américain (\sigle{GAO}) retrace l'attaque d'après les déclarations
de l'entreprise~\cite{gao2018equifax}.

\begin{tableaucours}{L{2.8cm}L{9.4cm}}{%
    \tableauHeader{Date (2017)} &
    \tableauHeader{Fait (d'après le rapport GAO-18-559)}
}
8 mars & L'organisme fédéral d'alerte en cybersécurité publie une
vulnérabilité du cadriciel Apache Struts, qui permet d'exécuter des
commandes à distance \\
10 mars & Des inconnus balaient les systèmes d'Equifax à la recherche de
cette vulnérabilité, et la trouvent sur le portail de contestation en ligne
des consommateurs \\
Mi-mars & Equifax diffuse l'alerte à ses administrateurs~; un balayage
interne, une semaine plus tard, ne détecte pas la vulnérabilité sur ce
portail \\
13 mai & Début de l'attaque proprement dite \\
13 mai -- 29 juillet & Les attaquants passent de 3 bases de données à 48
autres, lancent environ 9\,000 requêtes et extraient les données par
petites quantités \\
29 juillet & Un administrateur remplace un certificat expiré~; l'inspection
du trafic reprend et révèle l'intrusion, après environ 76 jours \\
\end{tableaucours}
\vspace{-0.8em}
\begin{center}{\small\itshape Tableau~2.3. La chronologie de l'attaque
d'Equifax (d'après~\cite{gao2018equifax})}\end{center}

Les données d'au moins 145,5~millions de personnes ont été
atteintes~\cite{gao2018equifax}. En juillet~2019, l'entreprise a conclu avec
les autorités fédérales et les États un accord pouvant atteindre
700~millions de dollars~\cite{cfpb2019equifax}.

\subsection{Les contrôles qui ont manqué}

L'entreprise a elle-même identifié, après enquête, quatre facteurs, plus un
cinquième~\cite{gao2018equifax}. Ils se traduisent en contrôles que tout
auditeur de sécurité sait nommer.

\begin{tableaucours}{L{2.8cm}L{5.0cm}L{4.4cm}}{%
    \tableauHeader{Facteur} &
    \tableauHeader{Ce qui s'est passé} &
    \tableauHeader{Contrôle défaillant}
}
Identification & La liste de diffusion de l'alerte était périmée~; ceux qui
devaient corriger ne l'ont pas reçue~; le balayage n'a pas vu la faille &
Gestion des vulnérabilités et des correctifs~; inventaire des actifs \\
Détection & Un certificat expiré empêchait l'outil d'inspecter le trafic
chiffré & Surveillance~; gestion des certificats \\
Segmentation & Les bases n'étaient pas isolées les unes des autres &
Cloisonnement du réseau et des données \\
Gouvernance des données & Une base contenait des identifiants en clair
donnant accès à d'autres bases & Gestion des secrets et des accès \\
Fréquence des requêtes & Aucune limite~: 9\,000 requêtes sans alarme &
Détection d'anomalies \\
\end{tableaucours}
\vspace{-0.8em}
\begin{center}{\small\itshape Tableau~2.4. Les facteurs de l'attaque
d'Equifax et les contrôles correspondants}\end{center}

\subsection{Un contrôle en place, qui ne fonctionnait pas}

Le facteur «~détection~» mérite une attention particulière. Equifax
\emph{avait} installé un outil d'inspection du trafic~: sur le papier, le
contrôle existait. Mais, selon le GAO, le certificat qui lui permettait de
lire le trafic chiffré avait expiré environ dix mois avant
l'attaque~\cite{gao2018equifax}. Pendant tout ce temps, l'outil ne voyait
rien, et personne ne s'en était aperçu.

C'est la distinction fondamentale entre la \emph{conception} d'un contrôle
(il existe, il est adapté au risque) et son \emph{efficacité de
fonctionnement} (il a réellement fonctionné sur la période). Un auditeur
qui aurait vérifié l'existence de l'outil aurait coté \OUI{}~; celui qui
aurait vérifié qu'il produisait des alertes sur la période aurait coté
\NON{}. Le Chapitre~\ref{chap:controle-interne} développe cette
distinction.

\begin{boiteRemarque}{Dix mois ou dix-neuf mois~?}
Le rapport du GAO indique que le certificat avait expiré environ dix mois
avant l'attaque~\cite{gao2018equifax}. Le rapport de la commission de
contrôle de la Chambre des représentants, publié en décembre~2018, indique
pour sa part que le dispositif de surveillance était inactif depuis
dix-neuf mois~\cite{house2018equifax}. Les deux rapports ne décrivent
peut-être pas le même équipement, ni le même point de départ. Un auditeur
ne tranche pas entre deux sources primaires sans pièce supplémentaire~: il
cite les deux, et dit que l'écart reste à expliquer. Ici, la conclusion
ne dépend pas de l'écart~: dans les deux cas, la surveillance était aveugle
pendant toute l'attaque.
\end{boiteRemarque}

\begin{boxmetier}{RSSI}
\textbf{Les questions que l'affaire lui pose}
\begin{enumerate}
    \item Ma liste de diffusion des alertes de sécurité est-elle à jour, et
          qui vérifie qu'un correctif critique a été installé partout~?
    \item Mon inventaire des applications exposées sur Internet est-il
          complet~? Le balayage couvre-t-il tout l'inventaire~?
    \item Qui suit la date d'expiration des certificats de mes outils de
          surveillance, et qui s'aperçoit qu'un outil ne voit plus rien~?
    \item Une base compromise donne-t-elle accès aux autres~?
\end{enumerate}
\end{boxmetier}

% ============================================================================
\section{Wirecard (2020) : la Trésorerie qui n'Existait pas}
\label{sec:c02-wirecard}
% ============================================================================

\subsection{Les faits}

Wirecard était une société allemande de traitement des paiements, entrée
en 2018 dans le principal indice boursier allemand. Une partie
importante de son activité passait par des partenaires tiers en Asie, et
les fonds correspondants étaient présentés comme déposés sur des comptes
fiduciaires.

\begin{tableaucours}{L{2.8cm}L{9.4cm}}{%
    \tableauHeader{Date} &
    \tableauHeader{Fait}
}
2015-2019 & Des articles de presse et des vendeurs à découvert mettent en
doute les comptes de la société \\
Février 2019 & L'autorité allemande de surveillance financière (BaFin)
interdit pour deux mois les ventes à découvert de l'action~; des plaintes
visent des journalistes et des vendeurs à découvert \\
Octobre 2019 & Sous la pression des investisseurs, la société confie une
enquête spéciale au cabinet KPMG \\
28 avril 2020 & Rapport de KPMG~: le cabinet n'a pas pu vérifier les
revenus tirés des partenaires tiers, ni une partie importante des soldes
de trésorerie~; il n'a reçu, presque toujours, que des copies
électroniques de documents dont il ne pouvait pas vérifier
l'authenticité~\cite{wirecard2020kpmg} \\
18 juin 2020 & Le commissaire aux comptes, EY, refuse de certifier les
comptes 2019~: il ne peut pas confirmer l'existence de 1,9~milliard
d'euros de soldes bancaires \\
22 juin 2020 & La société déclare qu'il est probable que ces soldes
n'existent pas \\
25 juin 2020 & Demande d'ouverture d'une procédure d'insolvabilité \\
\end{tableaucours}
\vspace{-0.8em}
\begin{center}{\small\itshape Tableau~2.5. La chute de Wirecard (d'après le
rapport de KPMG~\cite{wirecard2020kpmg} et les communiqués de la société,
rapportés par la presse)}\end{center}

\subsection{Le mécanisme et l'audit}

La trésorerie déclarée reposait sur des confirmations et des relevés
fournis par des intermédiaires, et non obtenus directement des banques.
Selon la presse, deux banques philippines dont les noms figuraient sur ces
documents ont déclaré n'avoir jamais eu de relation avec la société pour
ces comptes. Le commissaire aux comptes avait certifié les comptes
pendant plusieurs années~; c'est l'enquête spéciale, puis le refus de
certifier, qui ont fait apparaître que la preuve de l'existence de la
trésorerie n'avait jamais été établie de façon indépendante.

La règle d'audit en cause est ancienne~: la \termecle{confirmation
directe}\index{confirmation directe}. Pour s'assurer qu'un solde bancaire
existe, l'auditeur demande lui-même à la banque de le confirmer, et reçoit
la réponse directement, sans qu'elle passe par l'entité auditée. Un relevé
remis par l'audité, ou par un tiers qu'il désigne, ne vaut pas
confirmation.

\subsection{Ce qui a changé}

L'affaire a conduit à la démission des dirigeants de l'autorité de
surveillance financière, à une commission d'enquête du Bundestag ouverte en
2020, et à une réforme du contrôle de l'information financière et de la
supervision des auditeurs en Allemagne. Elle a aussi rappelé que le
régulateur lui-même peut se tromper de cible~: en 2019, il protégeait la
société contre ceux qui l'accusaient.

\begin{boxscenario}{Opération LIZBIZ. Un relevé de trop}
\textbf{Situation.} Pour justifier la trésorerie de son agence de
Bafoussam, LiZbiZ remet à AUDIX CONSEIL un relevé bancaire au format PDF,
envoyé par courriel par le comptable de l'agence.

\medskip
\textbf{Analyse} (confiance élevée). Le relevé est une copie électronique
remise par l'audité~: c'est exactement le type de pièce que KPMG n'a pas
pu authentifier chez Wirecard. Tant que la banque n'a pas confirmé le
solde directement à l'auditeur, la cote est \NEK{}, et non \OUI{}. La
confirmation est demandée par AUDIX CONSEIL, sur papier à en-tête de
LiZbiZ signé par le directeur financier, et la banque répond à l'adresse du
cabinet.
\end{boxscenario}

\begin{boxmetier}{Auditeur SI}
\textbf{Authentifier une pièce numérique}
\begin{enumerate}
    \item Préférer la pièce obtenue directement à la source (extraction
          réalisée en présence de l'auditeur, confirmation de la banque) à
          la pièce remise par l'audité.
    \item Pour un fichier remis, consigner sa provenance, sa date et son
          empreinte, et rechercher une corroboration indépendante.
    \item Coter \NEK{} tant que l'authenticité n'est pas établie~; ne
          jamais coter \OUI{} sur une copie seule.
\end{enumerate}
\end{boxmetier}

% ============================================================================
\section{Cinq Schémas qui Reviennent}
\label{sec:c02-schemas}
% ============================================================================

\subsection{La synthèse des quatre affaires}

\begin{tableaucours}{L{2.0cm}L{3.4cm}L{3.4cm}L{3.4cm}}{%
    \tableauHeader{Affaire} &
    \tableauHeader{Contrôle qui a manqué} &
    \tableauHeader{Ce que l'on a cru} &
    \tableauHeader{Ce qui a changé}
}
Enron & Surveillance du conflit d'intérêts~; information du conseil &
Que les montages étaient approuvés et compris & Loi Sarbanes-Oxley,
indépendance, rapport sur le contrôle interne \\
Société Générale & Fermeture des alertes sur preuve~; suivi en nominal &
Les explications du trader & Suivi du nominal, des annulations, des
congés~; supervision transversale \\
Equifax & Correctifs, inventaire, surveillance, segmentation & Qu'un outil
installé était un outil qui fonctionne & Accord de 2019, exigences de
sécurité renforcées \\
Wirecard & Confirmation directe de la trésorerie & Des copies de
documents remises par des tiers & Réforme de la supervision en
Allemagne \\
\end{tableaucours}
\vspace{-0.8em}
\begin{center}{\small\itshape Tableau~2.6. Quatre affaires lues avec la même
grille}\end{center}

\subsection{Cinq schémas}

Lues ensemble, ces affaires font apparaître cinq schémas de défaillance.
Ils se retrouvent dans les audits ordinaires, à une autre échelle.

\begin{enumerate}
    \item \textbf{Le contrôle écrit, non appliqué.} Il existe une règle, une
          procédure, un outil~; personne ne vérifie qu'il fonctionne
          (Equifax, Enron).
    \item \textbf{L'alerte fermée par une déclaration.} Le contrôle a
          réagi, mais une explication non prouvée a suffi à le faire taire
          (Société Générale).
    \item \textbf{La preuve fournie par l'intéressé.} L'auditeur accepte une
          pièce remise par la personne ou l'entité qu'il contrôle, sans
          confirmation indépendante (Wirecard).
    \item \textbf{L'indépendance compromise.} Celui qui évalue a intérêt à
          ne pas voir~: honoraires, proximité, relations (Enron).
    \item \textbf{La surveillance aveugle.} Le dispositif de détection
          existe mais ne voit plus rien, et personne ne surveille la
          surveillance (Equifax).
\end{enumerate}

\begin{boxinfo}
\textbf{Ce que dit la recherche.} Les schémas~2 et~3 illustrent l'axiome
d'\emph{interaction stratégique} de l'article \emph{Vers
l'auditologie}~\cite{minka2026auditologie}~: l'audité n'est pas une source
neutre, il a intérêt à la réponse. Les schémas~1 et~5 illustrent
l'\emph{observabilité partielle}~: l'auditeur ne voit le système qu'à
travers des traces, et une trace absente (un outil qui ne journalise plus)
ne dit pas que tout va bien. Dans les quatre affaires, ce qui aurait dû
être coté \NEK{} a été traité comme un \OUI{}.
\end{boxinfo}

\subsection{Les mêmes schémas chez LiZbiZ}

Les cinq schémas se retrouvent, à leur échelle, dans le cas fil rouge.
Les chapitres suivants les instruiront un à un.

\begin{tableaucours}{L{3.6cm}L{6.0cm}c}{%
    \tableauHeader{Schéma} &
    \tableauHeader{Chez LiZbiZ} &
    \tableauHeader{Chapitre}
}
Contrôle écrit, non appliqué & La procédure de désactivation des comptes
sous 48~heures~: 5 comptes d'agents partis encore actifs &
\ref{chap:preuve} \\
Alerte fermée par une déclaration & La première mission retient les
réponses majoritaires, sans preuve & \ref{chap:prologue} \\
Preuve fournie par l'intéressé & Les explications des ajustements de stock
de fin de mois, données par celui qui les saisit & \ref{chap:fraude} \\
Indépendance compromise & La première mission conduite par une équipe
dépendant de la direction évaluée & \ref{chap:audit-auditologie} \\
Surveillance aveugle & Aucun journal des revues de droits~; le
responsable informatique détient tous les droits de l'application &
\ref{chap:controle-interne} \\
\end{tableaucours}
\vspace{-0.8em}
\begin{center}{\small\itshape Tableau~2.7. Les cinq schémas dans le cas
LiZbiZ}\end{center}

\begin{boxmetier}{Gestionnaire de risques}
\textbf{Les risques que les affaires rendent visibles}

Aucune de ces affaires n'était imprévisible au sens strict~: le conflit
d'intérêts, la fraude interne, la faille non corrigée, la trésorerie non
confirmée sont des risques connus. Ce qui a manqué, c'est la vérification
que les mesures de traitement fonctionnaient. Le gestionnaire de risques
en retient qu'un risque n'est « traité » que lorsque l'efficacité de sa
mesure a été démontrée (Chapitre~\ref{chap:gestion-risques}).
\end{boxmetier}

\begin{boxmetier}{Responsable de la conformité}
\textbf{Le coût de la non-conformité}

L'accord conclu par Equifax en 2019 rappelle que la protection des données
personnelles a un coût, et que les autorités peuvent l'exiger a
posteriori. Au Cameroun, la loi n°~2024/017 sur les données personnelles
fait de la sécurité des traitements une obligation légale~; le
Chapitre~\ref{chap:audit-conformite} montre comment l'auditer.
\end{boxmetier}

\begin{boxinfo}
\textbf{Et les affaires africaines~?} Les défaillances survenues en
Afrique centrale sont rarement documentées par des sources publiques
complètes~: rapports d'enquête non publiés, procédures confidentielles. Le
cours ne cite que des affaires dont les sources primaires sont accessibles.
L'exercice~2.10 invite à en rechercher une, avec la même exigence.
\end{boxinfo}

\separateur

\begin{boxresume}
\textbf{Ce qu'il faut retenir de ce chapitre~:}
\begin{itemize}
    \item \textbf{Une affaire se lit avec une grille}~: faits, mécanisme,
          contrôle manquant, signal ignoré, défaut d'audit, réforme.
    \item \textbf{Les sources se critiquent}~: primaire ou secondaire,
          neutre ou partie~; une divergence se signale, ne se tranche pas en
          silence.
    \item \textbf{Enron} a conduit à la loi Sarbanes-Oxley~: sa section~404
          impose un rapport sur le contrôle interne, et a fait des
          contrôles informatiques un objet d'audit.
    \item \textbf{Société Générale}~: 75 alertes, toutes fermées par des
          explications~; une alerte ne se ferme que sur preuve.
    \item \textbf{Equifax}~: un correctif non installé, une surveillance
          aveugle depuis des mois~; un contrôle installé n'est pas un
          contrôle qui fonctionne.
    \item \textbf{Wirecard}~: une trésorerie jamais confirmée
          directement~; une copie remise par l'audité ne prouve pas.
    \item \textbf{Cinq schémas}~: contrôle écrit non appliqué, alerte
          fermée par une déclaration, preuve fournie par l'intéressé,
          indépendance compromise, surveillance aveugle. Tous reviennent à
          traiter un \NEK{} comme un \OUI{}.
\end{itemize}
\end{boxresume}

\begin{boxinfo}
\textbf{Points de contrôle.} Vous avez compris ce chapitre si vous savez~:
\begin{itemize}
    \item appliquer la grille des six questions à une affaire que vous ne
          connaissiez pas~;
    \item dire si une source est primaire ou secondaire, neutre ou partie~;
    \item expliquer, en trois phrases, pourquoi la section~404 a changé
          l'audit des systèmes d'information~;
    \item distinguer la conception d'un contrôle de son efficacité de
          fonctionnement, sur l'exemple d'Equifax~;
    \item reconnaître les cinq schémas dans une situation nouvelle.
\end{itemize}
\end{boxinfo}

% ============================================================================
\section*{Exercices}
% ============================================================================

\begin{boxexo}[title={\ding{45}\quad Exercice~2.1. Primaire ou secondaire~? \niveaubase}]
Classez chacune des sources suivantes (primaire, secondaire, partie) et
dites ce qu'elle permet d'établir.
\begin{enumerate}[(a)]
    \item Le rapport GAO-18-559 sur l'attaque d'Equifax.
    \item Un article d'encyclopédie en ligne sur Wirecard.
    \item Le document de la Société Générale sur l'affaire de 2008.
    \item Le texte de la section~404 de la loi Sarbanes-Oxley.
    \item Le communiqué par lequel Wirecard publie le rapport de KPMG.
\end{enumerate}
\end{boxexo}

\begin{boxexo}[title={\ding{45}\quad Exercice~2.2. La perte de la Société Générale \niveaubase}]
\begin{enumerate}[(a)]
    \item Retrouvez, à partir du Tableau~2.2, comment se forme la perte de
          4,9~milliards d'euros.
    \item Combien de fois les positions de janvier~2008 dépassaient-elles
          la limite de l'équipe~?
    \item Pourquoi une limite exprimée en position nette ne pouvait-elle
          pas révéler la fraude~? Quel indicateur l'aurait révélée~?
\end{enumerate}
\end{boxexo}

\begin{boxexo}[title={\ding{45}\quad Exercice~2.3. La chronologie d'Equifax \niveaumoyen}]
À partir du Tableau~2.3~:
\begin{enumerate}[(a)]
    \item Combien de jours se sont écoulés entre la publication de la
          vulnérabilité et le début de l'attaque proprement dite~? Qu'aurait
          dû permettre ce délai~?
    \item Pour chacun des cinq facteurs du Tableau~2.4, rédigez une question
          d'audit atomique, cotable \OUI{}, \NON{} ou \NEK{}, et la preuve
          attendue.
    \item Lequel de ces contrôles, s'il avait fonctionné, aurait le plus
          réduit la durée de l'attaque~? Justifiez.
\end{enumerate}
\end{boxexo}

\begin{boxexo}[title={\ding{45}\quad Exercice~2.4. Conception ou fonctionnement~? \niveaumoyen}]
Pour chacune des situations suivantes, dites si le défaut porte sur la
conception du contrôle ou sur son fonctionnement, et quel test l'aurait
révélé.
\begin{enumerate}[(a)]
    \item L'outil d'inspection du trafic d'Equifax, avec un certificat
          expiré.
    \item Une procédure de LiZbiZ qui prévoit la désactivation des comptes
          sous 48~heures, mais que les ressources humaines n'informent pas
          des départs.
    \item Une limite de risque exprimée en position nette, alors que les
          positions peuvent être compensées par des opérations fictives.
    \item Un comité d'audit qui reçoit des présentations trop générales
          pour comprendre les montages qu'on lui soumet.
\end{enumerate}
\end{boxexo}

\begin{boxexo}[title={\ding{45}\quad Exercice~2.5. Fermer une alerte \niveaumoyen}]
Le contrôle des paiements de LiZbiZ signale, sur l'exercice 2025,
vingt-trois paiements de 4,9~millions FCFA au même fournisseur, juste sous
le seuil de double signature de 5~millions. Mme~OWONA, qui les a saisis, explique
qu'il s'agit de livraisons fractionnées à la demande du fournisseur.
\begin{enumerate}[(a)]
    \item Comment l'alerte aurait-elle été traitée selon la « réponse 1 »
          de la Société Générale~? Selon la « réponse 3 »~?
    \item Quelles preuves indépendantes de Mme~OWONA faut-il obtenir avant
          de fermer l'alerte~?
    \item Quelle est la cote tant que ces preuves manquent~?
\end{enumerate}
\end{boxexo}

\begin{boxexo}[title={\ding{45}\quad Exercice~2.6. La confirmation directe \niveaumoyen}]
\begin{enumerate}[(a)]
    \item Décrivez, étape par étape, une procédure de confirmation directe
          d'un solde bancaire, en précisant qui envoie, qui reçoit et par
          quel canal.
    \item Pourquoi un relevé PDF remis par l'audité ne suffit-il pas, même
          s'il porte le logo de la banque~?
    \item Proposez deux autres situations, en audit des SI, où une
          confirmation directe auprès d'un tiers remplace une pièce remise
          par l'audité.
\end{enumerate}
\end{boxexo}

\begin{boxexo}[title={\ding{45}\quad Exercice~2.7. Section 404 et contrôles informatiques \niveaumoyen}]
Une société cotée utilise un progiciel de gestion intégré pour produire
ses états financiers.
\begin{enumerate}[(a)]
    \item Citez quatre contrôles informatiques dont le commissaire aux
          comptes devra évaluer l'efficacité pour attester le rapport de la
          direction.
    \item Pour chacun, indiquez la preuve que l'auditeur demandera.
    \item Pourquoi la section~404 a-t-elle accru la demande d'auditeurs des
          systèmes d'information~?
\end{enumerate}
\end{boxexo}

\begin{boxexo}[title={\ding{45}\quad Exercice~2.8. Les cinq schémas \niveaumoyen}]
Pour chacune des situations suivantes, identifiez le ou les schémas du
chapitre, et proposez le contrôle qui manque.
\begin{enumerate}[(a)]
    \item Le pare-feu de l'ANPP journalise les connexions, mais le disque
          des journaux est plein depuis quatre mois.
    \item L'auditeur interne de la BANQUE SANAGA évalue une procédure qu'il
          a rédigée.
    \item Le prestataire de LiZbiZ fournit lui-même le rapport mensuel sur
          sa disponibilité, sans autre source.
    \item Les écarts d'inventaire de l'agence de Garoua sont expliqués
          chaque mois par le chef d'agence, et l'explication est acceptée.
\end{enumerate}
\end{boxexo}

\begin{boxexo}[title={\ding{45}\quad Exercice~2.9. Analyser une affaire \niveauexpert}]
Choisissez l'une des affaires du chapitre et rédigez une note de deux pages
qui~:
\begin{enumerate}[(a)]
    \item applique les six questions du Tableau~2.1, en rattachant chaque
          fait à une source primaire~;
    \item rédige, pour le contrôle qui a manqué, une fiche d'assertion
          complète (objectif, critère, question, risque, conséquence), avec
          le test, la preuve et la trace~;
    \item dit ce qu'un auditeur, appliquant la règle « déclaration seule =
          \NEK{} », aurait vu, et à quel moment.
\end{enumerate}
\end{boxexo}

\begin{boxexo}[title={\ding{45}\quad Exercice~2.10. Une affaire africaine \niveauexpert}]
Recherchez une défaillance de contrôle interne ou de sécurité survenue en
Afrique, documentée par au moins une source primaire accessible (rapport
d'une juridiction des comptes, décision d'un régulateur, décision de
justice, rapport officiel).
\begin{enumerate}[(a)]
    \item Présentez la source et son statut (primaire, partie, neutre).
    \item Appliquez la grille des six questions, en distinguant ce que la
          source établit de ce qu'elle laisse ouvert (\NEK{}).
    \item Rattachez l'affaire à l'un des cinq schémas, ou montrez qu'elle en
          révèle un sixième.
\end{enumerate}
\end{boxexo}

\bigskip

\begin{boxinfo}
\textbf{Pour aller plus loin.}
\begin{itemize}
    \item Le rapport GAO-18-559~\cite{gao2018equifax}~: une quarantaine de
          pages, un modèle de rapport d'enquête factuel.
    \item Le rapport du comité spécial d'Enron~\cite{powers2002enron}, pour
          voir comment un conseil d'administration enquête sur lui-même.
    \item Le texte de la loi Sarbanes-Oxley~\cite{sox2002}, sections~201,
          302 et~404.
    \item Le rapport de KPMG sur Wirecard~\cite{wirecard2020kpmg}, pour voir
          ce qu'un auditeur écrit quand il ne peut pas conclure.
\end{itemize}
\end{boxinfo}

\begin{boxinfo}
\textbf{Transition.} Les quatre affaires ont un point commun~: un
dispositif de contrôle interne qui existait sur le papier et n'a pas
fonctionné. Le Chapitre~\ref{chap:controle-interne} définit ce dispositif,
ses composantes et ses principes, et montre comment le concevoir pour qu'il
fonctionne réellement.
\end{boxinfo}
